\properlane{communion-of-saints}{One people praises God across nations and death}

\subsection{The number John hears and the multitude he sees}
John's exact number can seem to restrict the hope proclaimed at All Saints.
He hears 144,000, then hears twelve thousand repeated for each named tribe.
The subsequent multitude exceeds counting and comes from every nation.
Bede addresses the relation between those descriptions at the definite
number itself. The number signifies the entire Church, whose members
descend from the patriarchs by birth or by imitation of faith. He supports
the inclusion of the latter with Paul's statement that those who belong to
Christ are Abraham's seed. His explanation concerns the way a people is
identified in the vision.\footnote{Bede, \work{Explanation}, I, on
Apoc.~7:4, p.~46; Gal.~3:29.}

The arithmetic is part of his argument. Twelve is already associated with
the apostles sent to preach the faith throughout the world. Its multiplication
by twelve and by thousands expresses the complete people gathered through
that mission. A precise number can therefore represent a multitude without
setting a statistical ceiling on it. Bede's explanation depends on the
symbolic force of the number and on membership by faith; it would lose its
warrant if the tribal list were treated merely as an unnecessary delay
before the multinational scene.

The tribal names receive extended attention in his exposition. Beginning
with Juda, which he associates with confession, Bede interprets the ordered
names through a sequence of virtues and their consummation in inheritance.
This is his spiritual exposition of the names. It does not require a reader
to accept the list as an independent lesson in Hebrew derivation.
Within his account, the named people are also a people formed in holiness.
The Gospel's descriptions of the blessed can consequently be compared with
the moral formation Bede finds in the enumeration, while each remains a
distinct biblical text.\footnote{Bede, \work{Explanation}, I, on
Apoc.~7:5--8, pp.~46--51.}

Bede's comment on verse 9 allows another account of the transition.
The enumeration can describe Israel, to whom the Gospel was first preached,
followed by the salvation of the Gentiles. This retains a difference between
the two groups at that point in the vision. It is an alternative to making
each scene describe precisely the same subjects in precisely the same way.
Both explanations admit the Gentiles to salvation; they locate the
distinction between Israel and the nations differently. The feast's
universal rejoicing can be understood through either, without deciding
that Bede supplied only one explanation.\footnote{Bede,
\work{Explanation}, I, on Apoc.~7:9, p.~52.}

\subsection{Belonging to the Lamb and confessing his gifts}
The seal identifies servants of the living God. The multitude's robes and
palms then describe people who stand in Christ's victory. Bede's baptismal
interpretation of the robes concerns their admission; his interpretation of
the palms concerns victory through the Cross. Salvation is acclaimed to
God and the Lamb, rather than claimed by the multitude as its own
achievement. Bede's comment on the sevenfold praise says that the Church
confesses, in each member, gifts received from God.\footnote{Bede,
\work{Explanation}, I, on Apoc.~7:9--10,12, pp.~52--53.}

This account helps explain the Introit's veneration of the saints alongside
its praise of the Son. The saints' holiness is a reason to praise the giver
of holiness. Schuster develops that relationship through the mystical body:
the Spirit's sanctification reaches its members in every condition.
His examples include people without social distinction as well as those
whose offices are conspicuous. The multitude's different languages can
accordingly remain visible. Their common confession of salvation does not
require that their earthly conditions or histories were identical.
\footnote{Schuster, \work{Sacramentary}, V, pp.~208--210.}

The angels' relation to this people also needs care. Bede permits them to
represent the worshipping multitude, but he also permits actual angelic
spirits rejoicing at the salvation of the Gentiles. The Introit explicitly
speaks of angels rejoicing and praising the Son. The Common Preface similarly
names heavenly powers and asks admission for human voices. Comparing these
prayers with Bede's latter explanation gives a straightforward account of
human and angelic praise together. His symbolic alternative remains an
explanation of the Apocalypse, without requiring that the Introit's
angels should mean human singers.\footnote{Bede, \work{Explanation}, I,
on Apoc.~7:11, p.~53; MR 1962, pp.~702, 252--253.}

The Preface's request is particularly exact: it asks that our voices be
admitted among theirs. Those praying are already giving thanks through
Christ, yet they ask to share in the heavenly acclamation. Their present
worship and the worship in John's vision are related through Christ and
directed toward God. Their conditions differ. That difference permits the
Mass to celebrate the saints' joy while retaining the congregation's
dependence on mercy.

\subsection{Death and continued membership in Christ}
Death raises a further question about this common people. The Offertory
sings of righteous souls whom the unwise suppose to have died, but who
remain in God's hand and peace. Augustine explains the same Wisdom phrase
while expounding Psalm 87. His difficulty there concerns the dead whom God
has not yet remembered by raising them. Righteous people are among the dead,
so their awaited resurrection cannot mean that they have been cut off from
God's protection. Wisdom 3:1 supplies his explicit reason: their souls are
in God's hands.\footnote{Augustine, Ps.~87, \S5; NPNF I.8,
Psalm LXXXVIII, \S4.}

Augustine distinguishes the Head from the members. Christ's death is freely
embraced; he has power to lay down his life and take it up. What has already
happened in his resurrection is reserved for his body at the end.
The righteous dead's continued security therefore concerns souls awaiting
resurrection. To say that they are at peace neither denies their bodily
death nor treats the final resurrection as already accomplished.
This makes the Offertory's apparent paradox intelligible: the persecutor can
kill the body without obtaining possession of the person as one separated
from God.

In \work{The City of God} XX.9, Augustine gives a further account of the
departed faithful's membership. His immediate subject is Apocalypse 20,
where souls reign with Christ. He appeals to remembrance of the departed at
the altar when the body of Christ is received. Such practice, he argues,
presupposes that the faithful dead remain Christ's members. Their souls
reign with him while they await their bodies. His argument concerns an
ecclesial relation that persists through death.\footnote{Augustine,
\work{The City of God}, XX.9, Dods edition, II (1871).}

The feast's prayers can be understood in that relation. Those celebrating
honor all the just, ask mercy through their intercession, and offer gifts
whose saving benefit they themselves still need. The saints' intercession
in the Collect and Postcommunion is expressly stated in those prayers.
Augustine's account of membership explains how death need not exclude the
departed from the Church's present communion; it does not supply his own
commentary on these particular orations. The connection with their
intercession follows from reading their actual petitions beside his
account of Christ's body.

\subsection{The present Church and the promised kingdom}
Augustine's account has an important distinction within it. The Church is
already Christ's kingdom, but its present condition includes wheat and
tares. The kingdom at the end excludes evil. The people gathered to
celebrate All Saints therefore cannot infer that everyone presently
within the assembly has already reached the holiness being celebrated.
The same chapter describes the saints' present reign through adherence to
Christ, while distinguishing it from the kingdom's final condition.
\footnote{Augustine, \work{City of God}, XX.9.}

The commemoration Collect makes this present need audible. The people
confess sins contracted through their frailty and ask release. The
feast's request for abundant mercy is thus joined to an account of why
mercy is needed. The Sunday Secret similarly asks God to complete an
unmerited gift and increase service. In this interpretation, the
congregation's membership becomes a reason for asking and serving.
The saints' communion encourages petition; it supplies no exemption
from conversion.

The Gradual and Gospel describe the life of those whom the Mass calls
saints. Fear of God, seeking him, mercy, purity, peace, and fidelity under
persecution are concrete relations and actions. Augustine's explanation
of upright praise already required submission to God. His explanation of
the Gradual names goods which do not depend on prosperity. The
Alleluia invites the burdened to Christ, so people who need pardon are
addressed within the celebration of those who have received his gifts.
The invitation concerns the present hearers and their need for the
Lord whom the saints praise.\footnote{Augustine, Ps.~32, enarratio I,
\S1; Ps.~33, enarratio II, \S\S14--15;
Chrysostom, \work{Matthew}, XXXVIII.3.}

Augustine's explanation of the Beatitudes describes their several rewards
as names of the one kingdom. Purity's sight of God, peace's divine sonship,
and the kingdom promised amid persecution are therefore compatible ends of
the same sanctified life. The Communion repeats precisely those three
promises. Those receiving hear them as promises still requiring fidelity,
while honoring saints whose happiness displays their fulfillment.
The Postcommunions continue that distinction: the faithful ask enduring
joy, the saints' prayers, and protection amid danger.\footnote{Augustine,
\work{Sermon on the Mount}, I.4.12; Matt.~5:8--10.}

\subsection{Four senses}
\begin{fourSenses}
\item[Literal.] John hears the enumeration of sealed Israel and sees a
multinational multitude praising God and the Lamb. The Mass honors all the
saints while its living worshippers ask mercy and protection.
\item[Allegorical.] Bede explains membership in Abraham's seed through
belonging to Christ. Augustine explains the continued membership of the
faithful dead in Christ's body. The Church gives thanks through Christ
with the heavenly company, while awaiting the body's resurrection.
\item[Moral.] Upright praise submits to God. Those celebrating the saints
confess frailty, seek God's good, practice mercy and peace, and remain
faithful under opposition. Present membership calls for the holiness
described in the Gospel.
\item[Anagogical.] The multitude's praise, the righteous dead's peace, and
the promised sight of God concern the completion sought by the pilgrim
Church. The saints' souls remain secure while the final resurrection
and the kingdom without evil are still awaited.
\end{fourSenses}
